克隆仓库并运行自动环境准备步骤：

\begin{lstlisting}[language=croqtile,caption={克隆仓库并准备依赖}]
git clone https://github.com/LancerLab/croqtile.git
cd croqtile
make setup-core
\end{lstlisting}

\co@make setup-core@ 会拉取 Git 子模块，并在需要时下载指定版本的 Flex、Bison、FileCheck 和 GoogleTest。环境准备完成后，构建编译器：

\begin{lstlisting}[language=croqtile,caption={构建编译器}]
make
\end{lstlisting}

接着运行测试套件，检查刚构建的编译器是否正常：

\begin{lstlisting}[language=croqtile,caption={运行测试套件}]
make test
\end{lstlisting}

所有测试通过后，就可以使用构建目录中的 \co@croq@ 可执行文件。将其所在目录加入 \co@PATH@，或者通过完整路径调用它。接下来，编译一个最小程序，验证整个安装流程。
